% ══════════════════════════════════════════════════════════════
\section{Conclusion}\label{sec:conclusion}
% ══════════════════════════════════════════════════════════════

Volatility forecasting in government bond markets is a problem of practical consequence and genuine difficulty. Bond volatility is non-stationary, regime-dependent, and driven by forces that shift in character over time: monetary policy stances, inflation dynamics, safe-haven flows, and central bank communication. Models that work well in one decade can deteriorate in the next, not because their structural logic is wrong but because the parameters that govern volatility dynamics and levels change with the economic environment. This paper asks whether these changes can be tracked in real time, without the overhead of re-estimation and without discarding the accumulated knowledge encoded in established volatility models.

The answer we provide is affirmative. The Model-Assisted Online Learning (MAOL) framework wraps any parametric volatility model in three adaptive layers, each with a distinct purpose and a distinct theoretical guarantee. Layer~1 updates the shape parameters of the volatility recursion observation by observation, allowing the model to track structural shifts in persistence and shock sensitivity as they occur. Layer~2 corrects slow-moving mis-scaling of the variance level through a multiplicative adjustment that responds directly to forecast errors without interfering with the dynamics estimated in Layer~1. Layer~3 recalibrates the predictive distribution in probability integral transform space, ensuring that the implied Value-at-Risk converges to the target exceedance rate regardless of the innovation distribution assumed. All three layers are driven by COCOB, a parameter-free optimizer that sets its own effective learning rate from the observed gradient history, eliminating the need for any manual tuning.

The empirical results speak directly to the framework's practical value. On 40 years of US Treasury futures, a purely sequential MAOL estimator with no access to future data achieves QLIKE within $0.00345$ units per step of a full-sample HAR-RV oracle that processes the entire sample at once. The difference is economically negligible and statistically insignificant by the Diebold-Mariano criterion. At the same time, MAOL-GARCH outperforms the EWMA industry benchmark by a margin that is both large (DM statistic $-34.5$) and robust: it holds without exception across all four bond markets tested. The variance-level multiplier $b_t$ provides a real-time, model-consistent signal of when the base recursion is adequate and when it is not. During the sharp COVID-19 volatility spike it adjusts within days, yielding a QLIKE advantage of 0.57 units per day over HAR Rolling. During the prolonged GFC it adapts more gradually, as theory predicts for a slow-moving structural shift. During tranquil periods it hovers near unity with negligible influence, confirming that the gradient signal drives adaptation only when genuine forecast errors accumulate.

Cross-country validation on German Bunds, UK Gilts, and Japanese JGBs, all using the same parameter settings without any market-specific adjustment, confirms that these properties are not specific to the US Treasury market. Average per-step regret against the within-market oracle is at most $+0.015$ QLIKE units in every market, VaR exceedance rates cluster near 5\%, and the DM tests reject equality with EWMA in every market at conventional significance levels. The Japanese JGB result is particularly telling: despite the Bank of Japan's yield curve control suppressing volatility for extended periods, the online calibration layer correctly identifies the altered tail properties and adjusts accordingly, without any model input describing monetary policy regime changes.

The sample-length study adds a dimension of robustness that is directly relevant to practitioners. Re-running MAOL-GARCH on non-overlapping windows of 2, 5, 10, and 15 years from the US Treasury history shows that forecast quality is broadly stable across window lengths. Even 2-year samples, roughly the shortest history available in an emerging fixed-income market, yield QLIKE values consistently below EWMA and VaR rates within the 3.9\% to 5.2\% band. The convergence phase occupies only the first one to two years of each window regardless of its length, meaning that a risk manager deploying MAOL in a market with limited history does not face a materially longer learning period than one with decades of data.

Together, these findings support a specific view of how online learning and econometric structure should interact. Established volatility models are not obstacles to be replaced by flexible algorithms: they encode empirical regularities that have survived across markets and decades, and discarding them in favour of black-box methods sacrifices interpretability without delivering commensurate gains in accuracy. The appropriate role for online learning is in the residual dynamics that parametric models leave unaddressed, whether because of structural change, unconditional level shifts, or distributional misspecification. MAOL operationalises this division of labour cleanly and with formal guarantees for each component.

Several directions remain open for future work. A multivariate extension incorporating online dynamic conditional correlation would enable joint covariance forecasting across fixed-income instruments, supporting portfolio-level risk aggregation. Alternative realized measures robust to microstructure noise and price jumps, such as bipower variation \citep{barndorff2004power} or the median realized variance estimator, could improve the quality of the proxy used for gradient computation. Incorporating scheduled macroeconomic events or central bank announcement indicators as additional covariates in Layer~1 could reduce forecast errors around high-impact information releases. Finally, the layer structure and COCOB optimizer apply without modification to richer base specifications such as asymmetric GARCH variants, long-memory HAR-type models, or semi-parametric volatility filters, all of which would inherit the same regret and calibration guarantees established here.
